\chapter{Acquisition modes}
\label{ch:acquisition}
\index{acquisition}

\section{Running a protocol}
\label{sec:running}

\begin{center}
\begin{tabularx}{\textwidth}{@{}lX@{}}
\toprule
Button & Behaviour \\
\midrule
\gui{Test Acquisition} & Runs the \emph{complete} protocol with all of the displays.
  Nothing is saved. Use the protocols with \file{\_test} in their names for a short
  check. \\
\gui{Start Acquisition} & The same, and each repetition is saved as it is recorded. \\
\gui{Pause} & Holds after the current repetition. \\
\gui{Resume} & Continues from where the protocol was paused. \\
\gui{Stop} & Abandons the protocol. The attenuators return to maximum. Files that were
  already written remain. \\
\bottomrule
\end{tabularx}
\end{center}

\index{Test Acquisition}\index{Start Acquisition}\index{Pause}\index{Stop}

Several acquisitions can be run in one session.\index{repeated acquisitions} While a
protocol is running or paused, \gui{Test Acquisition} and \gui{Start Acquisition} are
refused with a message. Let the protocol finish, or press \gui{Stop}, before starting
the next.

\section{Stimulus types}
\index{stimulus types}

The \key{stimulustype} entry in the protocol file selects the type of stimulus.

\begin{table}[htbp]
\centering
\caption{The protocols provided, with their approximate minimum running time if allowed
to finish. The time counts only the waveforms and the pause after each one.}
\label{tab:protocols}
\begin{tabular}{@{}lllrr@{}}
\toprule
Protocol file & Type & Levels (dB SPL) & Freqs & Minutes \\
\midrule
\file{clickIO} & click & 90 to 10, 17 levels & -- & 7 \\
\file{click\_test} & click & 94 & -- & 0.1 \\
\file{tonepip} & tonepip & 0 to 90, 10 levels & 9 & 18 \\
\file{tonepip\_test} & tonepip & 75, 94 & 1 & 0.4 \\
\file{interleaved\_plateau} & interleaved\_plateau & 10 to 90, 9 levels & 7 & 6 \\
\file{interleaved\_plateau\_high} & interleaved\_plateau & 60 to 90, 4 levels & 7 & 1.3 \\
\file{interleaved\_plateau\_low} & interleaved\_plateau & 10 to 50, 5 levels & 7 & $<$0.1 \\
\file{interleaved\_plateau\_test} & interleaved\_plateau & 30, 90 & 2 & 0.1 \\
\bottomrule
\end{tabular}
\end{table}

\begin{caution}
\file{interleaved\_plateau\_low.cfg} currently has \key{nreps} set to 5. The intended
value is 512. Correct the file before using this protocol.
\end{caution}

\begin{caution}
\file{interleaved\_ramp.cfg} and \file{interleaved\_random.cfg} appear in the list, but
these stimulus types are not in current use and are disabled in the waveform code. They
may be enabled again in the future. Selecting one and starting an acquisition produces
an ``Unrecognized wavetype'' error.
\index{interleaved ramp}\index{interleaved random}
\end{caution}

\subsection{Clicks}
\index{click}

Each sweep is one waveform, \key{wave\_duration} seconds long. It holds \key{nstim}
clicks, each lasting \key{stimulus\_duration}, spaced by \key{interval}. The first starts at
\key{delay}~+~\key{interval}. If \key{alternate} is true, the clicks alternate in
polarity.

One waveform is used for all levels. The level is set only by the attenuator:

\[
\mathrm{attenuation} = \texttt{click\_reference\_dbspl} - \mathrm{requested\ level}
\]

\key{dblist} is the list of levels, in the order presented. \key{freqlist} is empty.

\subsection{Tone pips}
\index{tone pip}

These are built like clicks, with \key{nstim} pips of \key{stimulus\_duration}, shaped by
\key{stimulus\_risefall}. The frequencies are in \key{freqlist}. All levels of one
frequency are presented before the next frequency begins.

The attenuation is calculated from the calibration, with linear interpolation between
calibrated frequencies:

\[
\mathrm{attenuation} = \mathrm{maxdB}(f) - \mathrm{requested\ level}
\]

Before anything is presented, the program checks that the calibration can reach every
requested level (Section~\ref{sec:level-check}).

\subsection{Interleaved plateau}
\label{sec:interleaved}
\index{interleaved plateau}

All levels and frequencies are placed in one waveform, so that every combination is
presented in every sweep, rather than in separate sweeps as for clicks and tone pips.
\key{nreps} is therefore the number of sweeps for the whole protocol. Each pip starts \key{stimulus\_period} after the last one,
beginning at \key{delay}. Within the waveform, all of the frequencies at the first level
come first, then all of the frequencies at the second level, and so on.

Each pip is scaled in the waveform according to the calibration. The whole waveform is
then scaled to 10~V at its largest pip, and one attenuator setting applies to the whole
sweep. The dynamic range of the output therefore limits the range of levels that can be
combined in a protocol.

\begin{caution}
\textbf{Recordings with several interleaved sets.} Recordings that store more than one
waveform set, such as those from December 2024, repeated \key{nreps} sweeps once for each
set until stopped. The code at the time probably played them at levels below those
requested, by the smallest attenuation stored in the file (about 2.9~dB for those
recordings).
\end{caution}

The entries in the protocol file are used differently for this type:

\begin{itemize}
  \item \key{stimulus\_period} is the spacing of the pips.
  \item \key{dblist} and \key{freqlist} are text, and may contain an expression such as
        \texttt{"np.linspace(10.0, 90.0, 9)"}.
  \item \key{wave\_duration}, \key{nstim} and \key{interval} are not used. The waveform
        lasts as long as the pips require.
  \item \key{alternate} has no effect: every sweep has the same polarity.
\end{itemize}

\section{Levels that cannot be reached}
\label{sec:level-check}
\index{Cannot reach}\index{negative attenuation}

A level above the maximum for its frequency would need a negative attenuation. For
tone pips and interleaved protocols the maximum comes from the calibration in use. The
maximum for clicks is the entry \key{click\_reference\_dbspl} in the configuration
file. When \gui{Test Acquisition} or
\gui{Start Acquisition} is pressed, the program checks every level in the protocol. If
any cannot be reached, a ``Cannot reach'' message lists them with the maximum for each
frequency, nothing is presented and no folder is created. Reduce the maximum level or
the range, and save the result as a new protocol.

\section{Protocol entries}
\label{sec:protocol-entries}
\index{protocol entries}

\begin{center}
\begin{tabularx}{\textwidth}{@{}lX@{}}
\toprule
Entry & Meaning \\
\midrule
\key{stimulustype} & \texttt{click}, \texttt{tonepip} or \texttt{interleaved\_plateau} \\
\key{wave\_duration} & Duration of one sweep, s (click and tone pip only) \\
\key{stimulus\_duration} & Duration of each click or pip, s \\
\key{stimulus\_risefall} & Rise and fall time of pips, s \\
\key{nstim} & Stimuli per sweep (click and tone pip only) \\
\key{delay} & Delay added to the stimulus start times, s \\
\key{interval} & Spacing of stimuli within a sweep, s (click and tone pip only) \\
\key{nreps} & Sweeps for each stimulus. For interleaved protocols, sweeps for the whole
  protocol \\
\key{stimulus\_period} & Pause after each sweep, s. For interleaved protocols it is also
  the spacing of the pips \\
\key{alternate} & Alternate the polarity within a sweep (click and tone pip only) \\
\key{freqlist} & Frequencies, Hz \\
\key{dblist} & Levels, dB SPL \\
\key{default\_frequency}, \key{default\_spl} & Used only if the lists are absent \\
\key{record\_dur}, \key{LPF} & Saved with the data, but have no effect on acquisition \\
\bottomrule
\end{tabularx}
\end{center}

\section{Recording}
\index{recording length}\index{sample rate}

The RP2.1 records at 24414.0625~Hz. The length of a record is that of the waveform,
rounded up to a whole multiple of 1024 samples. Nothing is filtered during acquisition,
beyond any filtering in the amplifier. Filtering for display is applied when viewing
(Chapter~\ref{ch:viewing}).

\subsection{The record must fit in the buffer of the circuit}
\label{sec:record-buffer}
\index{buffer}\index{SerStore}\index{recording does not fit}\index{wave\_duration}

The ABR is stored in the buffer \texttt{sampled\_wave} of the RP2.1 circuit, a SerStore
component. \emph{The size of this buffer is fixed when the circuit is designed.} It is
the \emph{Size} of the SerStore in RPvdsEx, in 32-bit words, and the program cannot make it
larger. (The program writes the number of samples to the tag \texttt{sampled\_wave\_n},
which is connected to the Size input in the circuit, but the memory of the buffer is
allocated when the circuit is built. The program therefore uses the size compiled into
the circuit, which it reads from the circuit.) A record that is longer than the buffer
cannot be stored. When this happened on the rig (a 1.2~s record, 29696 samples, with a
buffer of 26624 samples), the data from 1090.55~ms to the end of the record repeated the
beginning of the buffer: the length was right, but the data were not valid.

\paragraph{The check.} When \gui{Test Acquisition} or \gui{Start Acquisition} is pressed,
the program works out the record needed for the longest waveform in the protocol and
compares it with the size of the buffer, read from the circuit. If it does not fit, a
message ``The recording does not fit in the buffer of the RP2.1 circuit'' gives the
sizes and what to change, nothing is presented, and no folder is created. The same check
is made again when each sweep is armed.

\paragraph{The calculation.}
\begin{enumerate}
\item Samples $=$ waveform duration $\times$ sample rate (24414.0625~Hz), rounded up to a
  whole number.
\item Rounded up to a whole multiple of 1024. This is the record, and the number of
  samples the program asks for.
\item The record must be no more than the size of the buffer.
\end{enumerate}
The longest waveform that fits is therefore
$\lfloor \mathrm{size}/1024 \rfloor \times 1024 / \mathrm{sample\ rate}$.

\begin{center}
\begin{tabular}{@{}rrrr@{}}
\toprule
Waveform (s) & Samples & Record (samples) & Record (s) \\
\midrule
0.645 (interleaved plateau) & 15748 & 16384 $= 16 \times 1024$ & 0.671 \\
1.05 & 25635 & 26624 $= 26 \times 1024$ & 1.0905 \\
1.2 (clicks and tone pips as supplied) & 29297 & 29696 $= 29 \times 1024$ & 1.2163 \\
2 & 48829 & 49152 $= 48 \times 1024$ & 2.0133 \\
5 & 122071 & 122880 $= 120 \times 1024$ & 5.0332 \\
10 & 244141 & 244736 $= 239 \times 1024$ & 10.0244 \\
\bottomrule
\end{tabular}
\end{center}

\begin{center}
\begin{tabular}{@{}rrrr@{}}
\toprule
Buffer (samples) & 24414.0625~Hz & 48828.125~Hz & 97656.25~Hz \\
\midrule
26624 & 1.0905~s & 0.5453~s & 0.2726~s \\
256000 & 10.4858~s & 5.2429~s & 2.6214~s \\
\bottomrule
\end{tabular}
\end{center}
The second table is the longest waveform that fits, for the two sizes that the circuit has
had. It shows that doubling the sample rate halves the longest record. The ABR is recorded
at 24414.0625~Hz (set in \texttt{setup\_RP21} in \file{pystim3.py}).

\paragraph{When the record is too long.} There are two cures.
\begin{itemize}
\item Make the protocol's \key{wave\_duration} short enough, using the calculation
  above. A protocol whose last stimulus ends at 1~s needs a waveform of a little more
  than that, not 1.2~s.
\item Make the buffer larger in the circuit. Open the circuit file that is loaded
  (\file{abrs\_v3.rcx} in \file{tdt\_circuits}, named in \texttt{find\_hardware} in
  \file{pystim3.py}) in RPvdsEx, and set the \emph{Size} of the SerStore that holds
  \texttt{sampled\_wave} to at least the record that is needed, in samples. For one
  channel stored as 32-bit floating point numbers, a word is a sample, and takes 4
  bytes. For example, a Size of 256000 is 1,024,000 bytes (about 1~MB), and holds
  10.4858~s at 24414.0625~Hz. For a 5~s waveform the record is 122880 samples, so a Size
  of 122880 or more is needed; for a waveform at four times the sample rate, the Size
  must be four times larger for the same duration. If the SerStore holds several
  channels, or compressed data, the number of samples per channel is the number of
  words $\times$ the compression $/$ the number of channels; the program takes this into
  account when it reads the size.
\end{itemize}
The memory of the RP2.1 is limited, and the other buffers in the circuit use it as well.
RPvdsEx shows the memory use of the circuit, which must not be exceeded. After the
circuit is changed, the program reads the new size when it is next started.
